\chapter{A comunicação hormonal}\label{ch:g8:hormonal-communication}

A \omterm{def:g8:puberty:puberty}{puberdade} os apresentou
(\cref{def:g8:puberty:hormones}); a pílula os explora
(\cref{def:g8:contraception:pill}); a onda quente do susto entregou um
deles no meio de uma corrida
(\cref{ex:g8:nervous-communication:cases}). O serviço postal merece a
visita dele: as \omterm{def:g8:hormonal-communication:glands}{glândulas} que escrevem as cartas, o \omterm{prop:g4:journey-of-food:absorb}{sangue} que as
leva, o truque de endereçamento que torna precisa uma transmissão
aberta --- e as alças de \omterm{prop:g8:hormonal-communication:feedback}{retroalimentação} que mantêm toda a
correspondência em equilíbrio.

\section{Glândulas, hormônios, alvos}

\begin{definition}[As glândulas endócrinas]\label{def:g8:hormonal-communication:glands}
Uma \emph{glândula}\index{glândula} cujo produto é liberado \emph{no
\omterm{prop:g4:journey-of-food:absorb}{sangue}} (e não por um canal, como o da \omterm{def:g4:journey-of-food:tract}{saliva}) é uma glândula
\emph{endócrina}, e os produtos dela são os
\emph{hormônios}\index{hormônio}. As principais escritoras: a
\emph{glândula de comando} da base do \omterm{def:g5:brain-in-command:system}{cérebro} --- quem dá a partida na
\omterm{def:g8:puberty:puberty}{puberdade} e rege quase todas as outras; a \emph{tireoide}, no pescoço
--- o hormônio dela define o ritmo de funcionamento do corpo inteiro; as
\emph{suprarrenais}, como toucas sobre os \omterm{def:g7:removing-wastes:kidneys}{rins} --- o hormônio do alarme
dos sustos e das corridas; o \emph{pâncreas} --- o hormônio dele
administra a \omterm{prop:g7:organs-at-work:bill}{glicose} do \omterm{prop:g4:journey-of-food:absorb}{sangue}; e os \emph{\omterm{def:g8:reproductive-systems:female}{ovários} e os \omterm{def:g8:reproductive-systems:male}{testículos}} ---
os hormônios do ciclo e da reconstrução do
\cref{ch:g8:reproductive-systems}.
\end{definition}

\begin{proposition}[Endereçamento por receptor]\label{prop:g8:hormonal-communication:address}
O \omterm{prop:g4:journey-of-food:absorb}{sangue} entrega todo \omterm{def:g8:hormonal-communication:glands}{hormônio} em toda parte
(\cref{ex:g8:puberty:sites}) --- e, mesmo assim, cada um age só nos
\emph{órgãos-alvo}\index{órgão-alvo} dele: as \omterm{def:g5:first-look-at-cells:cell}{células} construídas com
pontos de acoplamento, os \emph{receptores}, para aquele mensageiro.
Transporte por transmissão aberta, endereçamento do lado de quem
recebe: uma carta que toda casa recebe e que só é legível onde a chave
encaixa. (O acoplamento da \omterm{def:g8:nervous-communication:synapse}{sinapse} da
\cref{def:g8:nervous-communication:synapse} é o mesmo truque, em
escala de nanômetros.)
\end{proposition}
\admitted

\begin{example}[A viagem de um mensageiro: o hormônio do alarme]\label{ex:g8:hormonal-communication:alarm}
Um quase acidente de bicicleta: as \omterm{def:g8:hormonal-communication:glands}{glândulas} suprarrenais despejam o
\omterm{def:g8:hormonal-communication:glands}{hormônio} do alarme delas no \omterm{prop:g4:journey-of-food:absorb}{sangue} e, em segundos --- uma volta da
circulação (\cref{ex:g4:heart-and-blood:round}) ---, todo receptor
equipado responde de uma vez: o ritmo do \omterm{def:g4:heart-and-blood:heart}{coração} dispara, as vias
aéreas se alargam, o \omterm{prop:g7:digestion-and-nutrients:absorption}{fígado} libera a \omterm{prop:g7:organs-at-work:bill}{glicose} depositada (o banco do
\cref{ex:g7:digestion-and-nutrients:breadtrace}, acionado), os \omterm{def:g4:heart-and-blood:vessels}{vasos}
dos \omterm{def:g3:bones-and-muscles:muscle}{músculos} se alargam e os do \omterm{def:g4:journey-of-food:tract}{tubo digestivo} se estreitam (o
redirecionamento da \cref{prop:g7:organs-at-work:supply}, agora
mantido por \omterm{def:g8:hormonal-communication:glands}{hormônio}) e a \omterm{def:g1:the-five-senses:sense}{pele} fica pálida e formigando. Uma carta,
seis respostas, um estado coordenado: pronto.
\end{example}

\section{Equilíbrio por retroalimentação}

\begin{proposition}[O princípio da retroalimentação]\label{prop:g8:hormonal-communication:feedback}
Os sistemas hormonais se regulam por
\emph{retroalimentação}\index{retroalimentação}: a \omterm{def:g8:hormonal-communication:glands}{glândula} que
escreve \emph{monitora o resultado} das próprias cartas e diminui o
ritmo quando o resultado é alcançado --- um termostato, e não uma
torneira. A vitrine é a \omterm{prop:g7:organs-at-work:bill}{glicose} do \omterm{prop:g4:journey-of-food:absorb}{sangue}: depois de uma refeição, a
própria \omterm{prop:g7:organs-at-work:bill}{glicose} que sobe dispara o \omterm{def:g8:hormonal-communication:glands}{hormônio} de estoque do pâncreas;
com a \omterm{prop:g7:organs-at-work:bill}{glicose} depositada, o nível cai, e a secreção diminui junto;
entre as refeições, um \omterm{def:g8:hormonal-communication:glands}{hormônio} parceiro libera o depósito no sentido
inverso. O nível segue por uma faixa estreita, dia e noite, sem
motorista consciente nenhum.
\end{proposition}
\admitted

\begin{example}[Quando a retroalimentação falha: o diabetes]\label{ex:g8:hormonal-communication:diabetes}
No \emph{diabetes}\index{diabetes}, o sinal do \omterm{def:g8:hormonal-communication:glands}{hormônio} de estoque
falha --- o pâncreas para de escrevê-lo, ou os receptores param de
responder. A \omterm{prop:g7:organs-at-work:bill}{glicose} se acumula no \omterm{prop:g4:journey-of-food:absorb}{sangue} depois de cada refeição; a
\omterm{def:g7:organs-at-work:recovery}{recuperação} do \omterm{def:g7:removing-wastes:kidneys}{rim} é sobrecarregada e aparece açúcar na \omterm{def:g7:removing-wastes:kidneys}{urina} --- o
sinal clássico do \cref{ex:g7:removing-wastes:thrift}, enfim
explicado. O tratamento é capítulo aplicado: doses medidas do \omterm{def:g8:hormonal-communication:glands}{hormônio}
que falta, no horário das refeições --- a carta escrita com seringa.
Milhões de pessoas vivem bem exatamente com esse arranjo.
\end{example}

\begin{remark}[Os dois sistemas em parceria]\label{rem:g8:hormonal-communication:partnership}
O telégrafo e o correio da
\cref{prop:g8:nervous-communication:compare} são uma administração só.
A \omterm{def:g8:hormonal-communication:glands}{glândula} de comando da base do \omterm{def:g5:brain-in-command:system}{cérebro} é a interface: \omterm{def:g5:brain-in-command:system}{sistema
nervoso} acima, \omterm{def:g8:hormonal-communication:glands}{hormônios} abaixo --- o susto do ciclista correu primeiro
pelo \omterm{def:g8:nervous-communication:neuron}{nervo} e foi sustentado por \omterm{def:g8:hormonal-communication:glands}{hormônio}; a \omterm{def:g8:puberty:puberty}{puberdade} corre devagar,
por \omterm{def:g8:hormonal-communication:glands}{hormônio}, sob um regente instalado no \omterm{def:g5:brain-in-command:system}{cérebro}; as ovelhas do
\cref{ex:g8:reproduction-and-environment:lambs} transformam a duração
do dia (uma percepção nervosa) em \omterm{def:g8:hormonal-communication:glands}{hormônios} de reprodução. Fios
rápidos onde a velocidade importa, correio lento onde o programa
precisa rodar por meses, e escritórios de tradução no meio.
\end{remark}

\section{Ler um sistema hormonal}

\begin{method}[As quatro perguntas de qualquer hormônio]\label{met:g8:hormonal-communication:read}
Para qualquer \omterm{def:g8:hormonal-communication:glands}{hormônio} que você encontrar, estabeleça:
\begin{enumerate}
  \item \emph{quem escreve}: que \omterm{def:g8:hormonal-communication:glands}{glândula}, para o \omterm{prop:g4:journey-of-food:absorb}{sangue}, disparada
        por quê?
  \item \emph{quem recebe}: que \omterm{prop:g8:hormonal-communication:address}{órgãos-alvo} carregam os receptores
        dele, e o que cada um faz ao receber?
  \item \emph{\omterm{prop:g8:hormonal-communication:feedback}{retroalimentação}}: que resultado se reporta de volta, e
        como o excesso de escrita é evitado?
  \item \emph{falha e uso}: o que acontece quando a carta falha --- e a
        medicina imita essa carta (a pílula, as doses de quem tem
        \omterm{ex:g8:hormonal-communication:diabetes}{diabetes})?
\end{enumerate}
\end{method}

\begin{example}[O método no marcador de ritmo da tireoide]\label{ex:g8:hormonal-communication:thyroid}
Quem escreve: a tireoide, regida pela \omterm{def:g8:hormonal-communication:glands}{glândula} de comando. Quem
recebe: quase toda \omterm{def:g5:first-look-at-cells:cell}{célula} --- o \omterm{def:g8:hormonal-communication:glands}{hormônio} define o ritmo do fogo lento
da \cref{prop:g7:what-is-respiration:power}. \omterm{prop:g8:hormonal-communication:feedback}{Retroalimentação}: a
\omterm{def:g8:hormonal-communication:glands}{glândula} de comando lê o nível e ajusta a regência dela. Falha: de
menos, e o motor inteiro fica em marcha lenta --- cansaço, frio, tudo
mais devagar; demais, e ele dispara. Uso: o \omterm{def:g8:hormonal-communication:glands}{hormônio} que falta é um
comprimido diário --- entre as cartas por correio mais antigas da
medicina.
\end{example}

\section{Exercícios}

\begin{exercise}[$\star$]\label{exo:g8:hormonal-communication:1}
O que faz uma \omterm{def:g8:hormonal-communication:glands}{glândula} ser endócrina? Diga quatro delas, com um papel
hormonal para cada.
\end{exercise}

\begin{exercise}[$\star$]\label{exo:g8:hormonal-communication:2}
O \omterm{prop:g4:journey-of-food:absorb}{sangue} entrega em toda parte; os \omterm{def:g8:hormonal-communication:glands}{hormônios} agem em algum lugar.
Resolva o paradoxo.
\end{exercise}

\begin{exercise}[$\star$]\label{exo:g8:hormonal-communication:3}
Liste quatro das seis respostas ao \omterm{def:g8:hormonal-communication:glands}{hormônio} do alarme, cada uma com o
capítulo que a atende.
\end{exercise}

\begin{exercise}[$\star$]\label{exo:g8:hormonal-communication:4}
Enuncie o princípio da \omterm{prop:g8:hormonal-communication:feedback}{retroalimentação} e a vitrine dele.
\end{exercise}

\begin{exercise}[$\star$]\label{exo:g8:hormonal-communication:5}
O que falha no \omterm{ex:g8:hormonal-communication:diabetes}{diabetes}, e o que reaparece no \omterm{def:g7:removing-wastes:kidneys}{rim} por causa disso?
\end{exercise}

\begin{exercise}[$\star$]\label{exo:g8:hormonal-communication:6}
Onde os \omterm{def:g5:brain-in-command:system}{sistemas nervoso} e hormonal fazem interface? Dê um exemplo de
revezamento.
\end{exercise}

\begin{exercise}[$\star\star$]\label{exo:g8:hormonal-communication:7}
Por que um termostato é a imagem certa para a \omterm{prop:g8:hormonal-communication:feedback}{retroalimentação}, e uma
torneira é a errada?
\end{exercise}

\begin{exercise}[$\star\star$]\label{exo:g8:hormonal-communication:8}
Rode o \cref{met:g8:hormonal-communication:read} no \omterm{def:g8:hormonal-communication:glands}{hormônio} de
estoque do pâncreas.
\end{exercise}

\begin{exercise}[$\star\star$]\label{exo:g8:hormonal-communication:9}
Explique a seringa de quem tem \omterm{ex:g8:hormonal-communication:diabetes}{diabetes} e a \omterm{def:g8:contraception:pill}{pílula anticoncepcional}
como o mesmo tipo de ato, pela pergunta 4 do método.
\end{exercise}

\begin{exercise}[$\star\star$]\label{exo:g8:hormonal-communication:10}
Atribua os sistemas, com motivos: o início do tremor de frio; o
crescimento de uma barba; a faixa estreita da \omterm{prop:g7:organs-at-work:bill}{glicose} do \omterm{prop:g4:journey-of-food:absorb}{sangue}; o
reflexo patelar.
\end{exercise}

\begin{exercise}[$\star\star$]\label{exo:g8:hormonal-communication:11}
As ovelhas do \cref{ex:g8:reproduction-and-environment:lambs}
transformam \omterm{def:g6:exploring-our-environment:conditions}{luz} em cordeiros. Trace o revezamento pelos escritórios da
\cref{rem:g8:hormonal-communication:partnership}.
\end{exercise}

\begin{exercise}[$\star\star\star$]\label{exo:g8:hormonal-communication:12}
``Os receptores transformam uma transmissão aberta numa rede de linhas
privadas.'' Defenda a frase e depois puxe mais: o que precisa mudar
numa \omterm{def:g5:first-look-at-cells:cell}{célula} para ela \emph{entrar} na plateia de um \omterm{def:g8:hormonal-communication:glands}{hormônio} --- e o
que isso sugere sobre como os alvos da \omterm{def:g8:puberty:puberty}{puberdade} passaram a estar
escutando?
\end{exercise}

\section{Problema: O diário da glicose}

\begin{problem}[{Problema de fim de semana --- um dia de açúcar no
sangue, lido como correspondência}]\label{pb:g8:hormonal-communication:1}
O registro (inventado) de um dia num monitor contínuo de \omterm{prop:g7:organs-at-work:bill}{glicose}:
estável durante a noite; café da manhã às 8 --- uma subida e depois uma
descida calma até as 10; esporte ao meio-dia --- uma queda breve, logo
contida; um almoço pulado --- estável a tarde inteira mesmo assim; um
susto às 17 (um quase acidente de bicicleta) --- uma subida rápida e
breve; jantar às 19 --- subida e descida de novo.

\medskip\noindent\textbf{Parte I --- A manhã.}
\begin{enumerate}
  \item A estabilidade da noite: que correspondência de duas cartas
        segurou a faixa enquanto o dono dormia?
  \item A subida e a descida do café da manhã: diga o disparador,
        quem escreve, a carta e os receptores que depositaram a sobra.
  \item Por que o fato de a descida \emph{parar} --- sem mergulhar
        abaixo da faixa --- é, em si, a assinatura da
        \omterm{prop:g8:hormonal-communication:feedback}{retroalimentação}?
  \item O nível das 10 é igual ao das 7. O que a imagem do termostato
        prevê sobre o percurso da secreção nesse intervalo?
\end{enumerate}

\medskip\noindent\textbf{Parte II --- A tarde.}
\begin{enumerate}[resume]
  \item A queda e a contenção do esporte: que carta acionou o banco
        do \omterm{prop:g7:digestion-and-nutrients:absorption}{fígado}, e a conta de que capítulo estava sendo gasta?
  \item O almoço pulado não mudou nada visível. O trabalho silencioso
        de que \omterm{def:g8:hormonal-communication:glands}{hormônio} parceiro atravessou a tarde?
  \item A subida rápida do susto: diga quem escreve e o objetivo desse
        açúcar repentino --- pronto para quê, segundo o
        \cref{ex:g8:hormonal-communication:alarm}?
  \item Por que a subida do susto diminuiu dentro de uma hora, ao
        contrário do acúmulo duradouro do \omterm{ex:g8:hormonal-communication:diabetes}{diabetes}?
\end{enumerate}

\medskip\noindent\textbf{Parte III --- A consulta.}
\begin{enumerate}[resume]
  \item O registro de uma pessoa com \omterm{ex:g8:hormonal-communication:diabetes}{diabetes} no mesmo dia seria
        diferente em toda refeição. Descreva a curva do café da manhã
        sem tratamento e onde apareceria o sinal do \omterm{def:g7:removing-wastes:kidneys}{rim}.
  \item A versão tratada: o que a dose medida antes de cada refeição
        imita, e por que o tamanho dela precisa combinar com o prato?
  \item A médica chama o registro saudável de ``uma conversa que você
        nunca ouve''. Liste quem falou no dia --- \omterm{def:g8:hormonal-communication:glands}{glândulas}, cartas,
        receptores --- numa tabela-resumo ou num parágrafo.
  \item Encerre com a moral do capítulo numa frase: que tipo de
        sistema mantém um valor dentro de uma faixa a vida inteira,
        sem motorista?
\end{enumerate}
\end{problem}
